% block: computational-applications-heading page:43 type:chapter
\clearpage
\section{\BT{Some Applications to Computational Science and Data Science}{计算科学和数据科学中的若干应用}}\label{chap:computational-applications}
\setcounter{footnote}{0}
% block: computational-applications-abstract page:43 type:ordinary-paragraph
\noindent\textbf{\BT{Abstract}{摘要}}\quad
\BT{Linear systems of the form $A\boldsymbol{x}=\boldsymbol{b}$ arise in a rich variety of applications such as computational science and data science. For computational science, physical phenomena are often described as partial differential equations. For solving partial differential equations, the finite-difference methods and the finite element method are widely used, leading to a problem of solving linear systems. Large-scale electronic structure calculation for condensed matter physics and lattice quantum chromodynamics for particle physics require solving large-scale linear systems. For data science, regression and classification are fundamental tasks that also require solving linear systems. Minimizing or maximizing functions is one of the most important optimization problems, which requires solving linear systems when Newton's method is used. We will see in this chapter how the linear systems arise in these applications by using simple and concrete examples.}{形如 $A\boldsymbol{x}=\boldsymbol{b}$ 的线性方程组出现在计算科学和数据科学等多种多样的应用中. 在计算科学中, 物理现象常用偏微分方程描述. 有限差分方法和有限元方法广泛用于求解偏微分方程, 由此产生求解线性方程组的问题. 凝聚态物理中的大规模电子结构计算和粒子物理中的格点量子色动力学都需要求解大规模线性方程组. 在数据科学中, 回归和分类是基本任务, 同样需要求解线性方程组. 求函数的极小值或极大值是最重要的优化问题之一, 使用 Newton 方法时需要求解线性方程组. 本章将通过简单而具体的例子, 说明这些应用中如何产生线性方程组.}
